% VI43-DETAILED-HINTS
% VI43-DETAILED-HINT-01
\paragraph{Exercise VI.43.1.}
\begin{hint}
% PEDAGOGY-ENRICHED-VI
Use the projective Jacobian criterion at geometric points. Use the cubic equation together with the chosen point or tangent line, then interpret intersections with multiplicity rather than as distinct points only.
\end{hint}

% VI43-DETAILED-HINT-02
\paragraph{Exercise VI.43.2.}
\begin{hint}
% PEDAGOGY-ENRICHED-VI
Differentiate and use the characteristic hypothesis. Use the cubic equation together with the chosen point or tangent line, then interpret intersections with multiplicity rather than as distinct points only.
\end{hint}

% VI43-DETAILED-HINT-03
\paragraph{Exercise VI.43.3.}
\begin{hint}
% PEDAGOGY-ENRICHED-VI
Distinct positive-degree components of a plane curve must meet. Use the cubic equation together with the chosen point or tangent line, then interpret intersections with multiplicity rather than as distinct points only.
\end{hint}

% VI43-DETAILED-HINT-04
\paragraph{Exercise VI.43.4.}
\begin{hint}
% PEDAGOGY-ENRICHED-VI
Apply Bézout to degrees \(1\) and \(3\). Use the cubic equation together with the chosen point or tangent line, then interpret intersections with multiplicity rather than as distinct points only.
\end{hint}

% VI43-DETAILED-HINT-05
\paragraph{Exercise VI.43.5.}
\begin{hint}
% PEDAGOGY-ENRICHED-VI
Use the rational function \(\ell/m\) restricted to the cubic. Use the cubic equation together with the chosen point or tangent line, then interpret intersections with multiplicity rather than as distinct points only.
\end{hint}

% VI43-DETAILED-HINT-06
\paragraph{Exercise VI.43.6.}
\begin{hint}
% PEDAGOGY-ENRICHED-VI
A general line section has three points counted with multiplicity. Use the cubic equation together with the chosen point or tangent line, then interpret intersections with multiplicity rather than as distinct points only.
\end{hint}

% VI43-DETAILED-HINT-07
\paragraph{Exercise VI.43.7.}
\begin{hint}
% PEDAGOGY-ENRICHED-VI
The tangent at a flex cuts the divisor \(3O\). Use the cubic equation together with the chosen point or tangent line, then interpret intersections with multiplicity rather than as distinct points only.
\end{hint}

% VI43-DETAILED-HINT-08
\paragraph{Exercise VI.43.8.}
\begin{hint}
% PEDAGOGY-ENRICHED-VI
Subtract \(3O\) from the collinear line-section relation. Use the cubic equation together with the chosen point or tangent line, then interpret intersections with multiplicity rather than as distinct points only.
\end{hint}

% VI43-DETAILED-HINT-09
\paragraph{Exercise VI.43.9.}
\begin{hint}
% PEDAGOGY-ENRICHED-VI
Cancel the common \(-O\) term and use degree-one rigidity. Use the cubic equation together with the chosen point or tangent line, then interpret intersections with multiplicity rather than as distinct points only.
\end{hint}

% VI43-DETAILED-HINT-10
\paragraph{Exercise VI.43.10.}
\begin{hint}
% PEDAGOGY-ENRICHED-VI
Subtract the two line-section relations in the two-line construction. Use the cubic equation together with the chosen point or tangent line, then interpret intersections with multiplicity rather than as distinct points only.
\end{hint}

% VI43-DETAILED-HINT-11
\paragraph{Exercise VI.43.11.}
\begin{hint}
% PEDAGOGY-ENRICHED-VI
Transport the identity element through the Abel--Picard bijection. Use the cubic equation together with the chosen point or tangent line, then interpret intersections with multiplicity rather than as distinct points only.
\end{hint}

% VI43-DETAILED-HINT-12
\paragraph{Exercise VI.43.12.}
\begin{hint}
% PEDAGOGY-ENRICHED-VI
Transport commutativity through addition in \(\mathrm{Pic}^0(E)\). Use the cubic equation together with the chosen point or tangent line, then interpret intersections with multiplicity rather than as distinct points only.
\end{hint}

% VI43-DETAILED-HINT-13
\paragraph{Exercise VI.43.13.}
\begin{hint}
% PEDAGOGY-ENRICHED-VI
Transport associativity through addition in \(\mathrm{Pic}^0(E)\). Use the cubic equation together with the chosen point or tangent line, then interpret intersections with multiplicity rather than as distinct points only.
\end{hint}

% VI43-DETAILED-HINT-14
\paragraph{Exercise VI.43.14.}
\begin{hint}
% PEDAGOGY-ENRICHED-VI
Use the line through \(P\) and \(O\), then the relation with \(3O\). Use the cubic equation together with the chosen point or tangent line, then interpret intersections with multiplicity rather than as distinct points only.
\end{hint}

% VI43-DETAILED-HINT-15
\paragraph{Exercise VI.43.15.}
\begin{hint}
% PEDAGOGY-ENRICHED-VI
Use the vertical line in short Weierstrass form. Use the cubic equation together with the chosen point or tangent line, then interpret intersections with multiplicity rather than as distinct points only.
\end{hint}

% VI43-DETAILED-HINT-16
\paragraph{Exercise VI.43.16.}
\begin{hint}
% PEDAGOGY-ENRICHED-VI
Substitute the secant equation and compare the sum of the three roots. Use the cubic equation together with the chosen point or tangent line, then interpret intersections with multiplicity rather than as distinct points only.
\end{hint}

% VI43-DETAILED-HINT-17
\paragraph{Exercise VI.43.17.}
\begin{hint}
% PEDAGOGY-ENRICHED-VI
Implicitly differentiate the Weierstrass equation. Use the cubic equation together with the chosen point or tangent line, then interpret intersections with multiplicity rather than as distinct points only.
\end{hint}

% VI43-DETAILED-HINT-18
\paragraph{Exercise VI.43.18.}
\begin{hint}
% PEDAGOGY-ENRICHED-VI
Order two means \(P=-P\). Use the cubic equation together with the chosen point or tangent line, then interpret intersections with multiplicity rather than as distinct points only.
\end{hint}

% VI43-DETAILED-HINT-19
\paragraph{Exercise VI.43.19.}
\begin{hint}
% PEDAGOGY-ENRICHED-VI
A flex tangent cuts \(3P\), while the chosen flex origin gives \(H\sim3O\). Use the cubic equation together with the chosen point or tangent line, then interpret intersections with multiplicity rather than as distinct points only.
\end{hint}

% VI43-DETAILED-HINT-20
\paragraph{Exercise VI.43.20.}
\begin{hint}
% PEDAGOGY-ENRICHED-VI
A singular affine point is a repeated root of \(x^3+Ax+B\). Use the cubic equation together with the chosen point or tangent line, then interpret intersections with multiplicity rather than as distinct points only.
\end{hint}

% VI43-DETAILED-HINT-21
\paragraph{Exercise VI.43.21.}
\begin{hint}
% PEDAGOGY-ENRICHED-VI
Restrict the cubic equation to the \(k\)-rational secant or tangent line. Use the cubic equation together with the chosen point or tangent line, then interpret intersections with multiplicity rather than as distinct points only.
\end{hint}

% VI43-DETAILED-HINT-22
\paragraph{Exercise VI.43.22.}
\begin{hint}
% PEDAGOGY-ENRICHED-VI
Normalization separates branches: a node has two and a cusp has one. Use the cubic equation together with the chosen point or tangent line, then interpret intersections with multiplicity rather than as distinct points only.
\end{hint}

% VI43-DETAILED-HINT-23
\paragraph{Exercise VI.43.23.}
\begin{hint}
% PEDAGOGY-ENRICHED-VI
Move the two normalization points to \(0\) and \(\infty\). Use the cubic equation together with the chosen point or tangent line, then interpret intersections with multiplicity rather than as distinct points only.
\end{hint}

% VI43-DETAILED-HINT-24
\paragraph{Exercise VI.43.24.}
\begin{hint}
% PEDAGOGY-ENRICHED-VI
Move the single removed point to infinity. Use the cubic equation together with the chosen point or tangent line, then interpret intersections with multiplicity rather than as distinct points only.
\end{hint}
